% Einheit 1 von 3 – Koordinatensystem (bank/symmetrie-abbildungen/e1.jsonl, zusammenbau v0.1)
%% TODO Zweigzeile Teil 1: was hier gelernt wird (Satz in Schülersprache aus dem Einheitstitel) – Einheit 1
\einheitenkopf[e1]{Einheit 1 von 3 · Koordinatensystem}
\zweigzeile{ab Kl. 5 · P10}
\setcounter{aufgabe}{9}

%% TODO Titel: Ich-kann-Satz fehlt in der Bank (Platzhalter: Kettenname) – Nr. 10
%% TODO Anweisung: ein Satz über der ersten Teilaufgabe fehlt in der Bank – Nr. 10
\begin{aufgabe}{„Erst rechts, dann hoch“}
\swa{Zeichne vom Ursprung aus nur den Weg zu $(4|1)$ mit zwei Pfeilen ein: erst nach rechts, dann nach oben. Setze keinen Punkt.}{\begin{ksys}[xmin=0,xmax=7,ymin=0,ymax=7] \end{ksys}}
\end{aufgabe}

%% TODO Titel: Ich-kann-Satz fehlt in der Bank (Platzhalter: Kettenname) – Nr. 11
%% TODO Anweisung: ein Satz über der ersten Teilaufgabe fehlt in der Bank – Nr. 11
\begin{aufgabe}{Punkte und Figuren im Koordinatensystem}
\begin{teile}
\teil Du sollst $(5|1)$ eintragen. Wohin gehst du vom Ursprung zuerst? \\ \kreuz{zuerst nach rechts} \\ \kreuz{zuerst nach oben}
\end{teile}
\swa{Trage den Punkt $A(4|6)$ ein.}{\begin{ksys}[xmin=0,xmax=8,ymin=0,ymax=8] \end{ksys}}
\swa{Trage den Punkt $B(1|3)$ ein.}{\begin{ksys}[xmin=0,xmax=8,ymin=0,ymax=8] \end{ksys}}
\swa{Trage den Punkt $C(6|1)$ ein.}{\begin{ksys}[xmin=0,xmax=8,ymin=0,ymax=8] \end{ksys}}
\swa{Trage den Punkt $D(5|5)$ ein.}{\begin{ksys}[xmin=0,xmax=8,ymin=0,ymax=8] \end{ksys}}
\swa{Trage den Punkt $E(2|7)$ ein.}{\begin{ksys}[xmin=0,xmax=8,ymin=0,ymax=8] \end{ksys}}
\swfrage{Was bleibt gleich, was ändert sich?}
\swz[3]{Lies die Koordinaten des Punktes F ab.}{\begin{ksys}[xmin=0,xmax=8,ymin=0,ymax=8] \punkt{6}{4}{F} \end{ksys}}
\swa{Trage die Punkte $B(-4|1)$ und $C(2|-3)$ ein.}{\begin{ksys}[xmin=-6,xmax=6,ymin=-6,ymax=6] \end{ksys}}
%% TODO Darstellung für schwach fehlt in der Bank (symmetrie-abbildungen-e1-k2-s4-v1; Abschnitt „Für schwache Schüler“)
\swz[3]{In welchem Quadranten liegt $H(-3|4)$? Die Quadranten werden gegen den Uhrzeigersinn gezählt, rechts oben ist der 1.}{}
\swa{Trage $K(0|-4)$ ein. Auf welcher Achse liegt K?}{\begin{ksys}[xmin=-6,xmax=6,ymin=-6,ymax=6] \end{ksys}}
\swa{Trage $A(1|1)$, $B(6|1)$, $C(6|4)$, $D(1|4)$ ein und verbinde sie der Reihe nach. Wie heißt die Figur?}{\begin{ksys}[xmin=0,xmax=8,ymin=0,ymax=8] \end{ksys}}
\swa{$A(1|2)$, $B(6|2)$ und $C(6|5)$ sind drei Ecken des Rechtecks ABCD. Zeichne und gib die Koordinaten von D an. D( \leerfeld | \leerfeld )}{\begin{ksys}[xmin=-5,xmax=7,ymin=-5,ymax=7] \end{ksys}}
\begin{teile}
\teil Genau einer der vier Punkte liegt auf der Rechtsachse (x-Achse). Welcher? \hfill (P10 2020 OS) \\ \kreuz{$A(6|4)$} \\ \kreuz{$B(0|4)$} \\ \kreuz{$C(-6|-4)$} \\ \kreuz{$D(6|0)$}
\end{teile}
\end{aufgabe}

%% TODO Titel: Ich-kann-Satz fehlt in der Bank (Platzhalter: Kettenname) – Nr. 12
%% TODO Anweisung: ein Satz über der ersten Teilaufgabe fehlt in der Bank – Nr. 12
\begin{aufgabe}{Achsen und Ursprung benennen, Einteilung ablesen}
\swz[3]{Lies ab: Wie groß ist ein Schritt der Einteilung auf der Rechtsachse?}{\begin{ksys}[xmin=0,xmax=10,ymin=0,ymax=10,xstep=2,ystep=2] \end{ksys}}
\end{aufgabe}

%% TODO Titel: Ich-kann-Satz fehlt in der Bank (Platzhalter: Kettenname) – Nr. 13
%% TODO Anweisung: ein Satz über der ersten Teilaufgabe fehlt in der Bank – Nr. 13
\begin{aufgabe}{Schreibweise P(x | y) lesen und schreiben (erste Zahl nach rechts, zweite nach oben)}
%% TODO Darstellung für schwach fehlt in der Bank (symmetrie-abbildungen-e1-k4-s1-v1; Abschnitt „Für schwache Schüler“)
\swz[3]{Der Punkt Q liegt vom Ursprung aus 6 nach rechts und 1 nach oben. Schreibe Q mit Koordinaten.}{}
\end{aufgabe}

%% TODO Titel: Ich-kann-Satz fehlt in der Bank (Platzhalter: Kettenname) – Nr. 14
%% TODO Anweisung: ein Satz über der ersten Teilaufgabe fehlt in der Bank – Nr. 14
\begin{aufgabe}{Streckenlänge parallel zu einer Achse abzählen}
%% TODO Darstellung für schwach fehlt in der Bank (symmetrie-abbildungen-e1-k5-s1-v1; Abschnitt „Für schwache Schüler“)
\swz[3]{Wie lang ist die Strecke von $A(1|2)$ bis $B(7|2)$ in Kästchen?}{}
\end{aufgabe}

%% TODO Titel: Ich-kann-Satz fehlt in der Bank (Platzhalter: Kettenname) – Nr. 15
%% TODO Anweisung: ein Satz über der ersten Teilaufgabe fehlt in der Bank – Nr. 15
\begin{aufgabe}{Punkte und Figuren im Koordinatensystem – Fehler finden, Begründen, Darstellungswechsel, Anwendung}
\swz[3]{Finde den Fehler. Lea soll $A(2|6)$ eintragen. Sie geht vom Ursprung 6 nach rechts und dann 2 nach oben.}{}
\swfrage{Begründe: Liegt ein Punkt auf der Rechtsachse, ist seine zweite Koordinate 0.}
\swz[3]{Schreibe die drei Punkte aus dem Bild mit Koordinaten auf.}{\begin{ksys}[xmin=-3,xmax=7,ymin=-3,ymax=7] \punkt{1}{4}{A} \punkt{5}{6}{B} \punkt{6}{-2}{C} \end{ksys}}
\swz[3]{Auf dem Lageplan des Schulhofs ist der Eingang der Ursprung, ein Kästchen sind 5\,m. Die Tischtennisplatte liegt bei $T(8|6)$. Wie viele Meter liegt sie rechts vom Eingang?}{}
\end{aufgabe}

\uebersichtskasten{Koordinaten: Die beiden Achsen treffen sich im Ursprung. Zu jedem Punkt gehören zwei Zahlen: die erste sagt, wie weit nach rechts, die zweite, wie weit nach oben. Geschrieben wird P(3 | 2). \\ Vier Quadranten: Nach links und nach unten zählst du mit negativen Zahlen. Q(−3 | −2). \\ Auf einer Achse: Liegt ein Punkt auf der Rechtsachse, ist seine zweite Koordinate 0 – R(3 | 0). Liegt er auf der Hochachse, ist die erste Koordinate 0 – S(0 | 2). Im Ursprung sind beide 0.}
